\par\Needspace{7\baselineskip}\bigskip
\begin{center}
\pdfbookmark[2]{Gauss to Bessel, Göttingen, 21 March 1843}{pallas-brief-18430321}
\textsc{Gauss} to \textsc{Bessel}, Göttingen, 21 March 1843.
\end{center}
\noindent\hspace*{\parindent}\mbox{. . . . .}\quad The first paper [by \textsc{Bessel}] on Jupiter's mass stirred in me a memory that is always painful to me, namely of my old work on the perturbations of Pallas. It has grown so unfamiliar to me over nearly a quarter of a century that I find it difficult even to find my way through the surviving papers. Of course, the need to increase \textsc{Laplace}'s value for Jupiter's mass was recognized at once. I initially set it at \(\frac{48}{47}\), although I cannot at this moment discover when I made that determination. In a later calculation, however, I find that the elements were thoroughly improved on the basis of the first 10 oppositions, \(-2376\) days . . . . . \(+1759{,}5\) days (reckoned from the beginning of 1810 as the epoch), and that in doing so, inter\,alia, that fraction was also increased to \(\frac{43}{42}\). The sum of the squares of the remaining errors was found to be (NB. without taking other perturbations into account, in particular those due to Saturn):
\begin{center}
\renewcommand{\arraystretch}{1.3}
\setlength{\tabcolsep}{5pt}
\begin{tabular}{@{}lrr@{}}
\textsc{Laplace}'s mass \(\left[\frac{1}{1067{,}09}\right]\) & (Hel.) longitude \(3175{,}69\) & (Hel.) latitude \(348{,}54\) \\
increased by \(\frac{1}{47}\) & \(963{,}78\) & \(318{,}47\) \\
increased by \(\frac{1}{42}\) & \(918{,}46\) & \(322{,}66\).
\end{tabular}
\end{center}
\par\SourceLeaf{444}There is then another calculation to improve all the elements on the basis of the first 12 oppositions, but in the papers I can now find the calculation has not been carried through to the end. While searching, however, I found some older papers containing a similar calculation from the first 9 oppositions\unskip\nobreak\HistoricalNoteEight{p. 561 of this volume.]}, though this still does not allow me to say with certainty whether the correction to the mass was found earlier still. At the same time I see that what I wrote above from memory was wrong. The perturbations due to Saturn were indeed taken into account. It is probably superfluous to remark that all these calculations are based on the theory of perturbations that I call the general theory, since it gives the amount of the perturbations for any time whatever. In all there are more than 800 equations (if those that contain the sine and cosine of the same angle are combined, hence more than 1600 without this combination), and this theory had been worked through in full more than once, as was necessary in order first to obtain the mean elements themselves. Alongside this general theory the calculation by quadratures was continually carried on, and it provided an absolutely indispensable check on the theory. --- The last thing I find is the calculation and comparison of the 14\textsuperscript{th} opposition of 6 January 1820. Whether or when I shall get around to studying this work again and editing and publishing it remains to be seen. Without a change in my circumstances that would leave me entirely free to dispose of my time, it will hardly happen, particularly since so many other things still present themselves that I am just as reluctant to let go to waste, or more accurately, to which I myself attach greater value.
\par\bigskip\centerline{\rule{3cm}{.4pt}}\par
